\begin{textsample}{NEG Dim 2 – 2014\_c4 – Score 9.00 – 2014\_c4/enade\_2014\_eng\_contr\_b460d1a0ee.txt}  \label{ex:f2_neg_316}
CAPÍTULO 4 PERCEPÇÃO DA PROVA As análises feitas neste capítulo tratam das percepções dos concluintes da Área de Engenharia de Controle e Automação sobre a prova aplicada no ENADE/2014.
Estas percepções foram mensuradas por meio de nove questões que avaliaram desde o grau de dificuldade da prova até o tempo gasto para concluí -la.
As percepções sobre a prova foram relacionadas com o desempenho dos estudantes e com a Grande Região de funcionamento do curso.
O questionário de percepção da prova encontra -se ao final do Anexo IV, que traz a reprodução da prova.
O desempenho dos estudantes foi classificado em quatro quartos.
Para tanto, esse desempenho foi ordenado de forma ascendente.
O percentil 25, P25, também conhecido como primeiro quartil, é a nota de desempenho que deixa um quarto ( 25 \% ) dos valores observados abaixo e três quartos acima.
A Figura 1 apresenta uma ilustração deste conceito.
O quarto inferior de desempenho é composto pelas notas abaixo do primeiro quartil.
Já o percentil 75, P75, também conhecido como terceiro quartil, é o valor para o qual há três quartos ( 75 \% ) dos dados abaixo e um quarto acima dele.
O quarto superior de desempenho é composto pelas notas iguais ou acima do terceiro quartil.
O percentil 50, P50, também conhecido como mediana, é o valor que divide as notas em dois conjuntos de igual tamanho.
O segundo quarto inclui valores entre o primeiro quartil ( P25 ) e a mediana.
O terceiro quarto contém os valores entre a mediana ( P50 ) e o terceiro quartil ( P75 ).
Vale ressaltar que percentis, quartis e medianas são pontos que não obrigatoriamente pertencem ao conjunto original de dados, ao passo que os quartos são subconjuntos dos dados originais.
A seguir, serão apresentados gráficos com resultados selecionados, relativos às nove questões avaliadas por grupos de estudantes.
Os gráficos apresentam nas barras o percentual de alunos que assinalaram uma das opções ou a soma das porcentagens daqueles que assinalaram duas ( ou três ) delas.
Por exemplo, para as questões 1 e 2, os gráficos apresentam a porcentagem total de participantes que assinalaram as opções ( D ) difícil e ( E ) muito difícil.
Em cada barra foram assinalados também os extremos do intervalo de confiança de 95 \% como linhas verticais unidas por uma linha horizontal na forma da letra H maiúscula.
As Tabelas no Anexo II apresentam os valores absolutos e a distribuição percentual20 das alternativas válidas das nove questões, segundo o mesmo recorte de desempenho dos alunos e Grande Região de funcionamento do curso. 4.1 GRAU DE DIFICULDADE DA PROVA 4.1.1 Componente de Formação Geral Ao avaliarem “ Qual o grau de dificuldade desta prova na parte de Formação Geral? ” ( Questão 1 ), 23,2 \% do grupo de inscritos e presentes optaram pelas alternativas difícil ou muito difícil.
Entretanto, para mais da metade dos estudantes ( 58,3 \% ), o Componente de Formação Geral da prova foi considerado com grau de dificuldade médio ( Gráfico 4.1, Gráfico 4.2 e, no Anexo II, a Tabela II.1 ).
O percentual de estudantes que consideraram a prova como difícil ou muito difícil foi maior na região Sudeste, onde a proporção foi de 25,5 \%, enquanto a de menor incidência foi a Nordeste, com 14,2 \%.
No Gráfico 4.1, é possível observar que esta diferença é estatisticamente significativa.
Nas Grandes Regiões, a proporção de presentes à prova que consideraram o Componente de Formação Geral como sendo de grau de dificuldade médio esteve entre 70,7 \% na região Norte e 57,3 \% na região Sudeste.
O percentual de alunos que consideraram a prova difícil ou muito difícil foi decrescente em função dos quartos de desempenho : 35,2 \% no primeiro quarto e 10,5 \% no quarto, grupo de melhor desempenho na prova.
Nos quartos de desempenho intermediários, a proporção de alunos que consideraram a prova difícil ou muito difícil foi de 28,0 \% no 2º quarto e 19,2 \% no 3º quarto.
As diferenças entre os quartos são estatisticamente significativas.
Para todos os quartos de desempenho, a alternativa modal para esta pergunta foi médio, com 54,0 \% e 58,7 \% dos respondentes nos quartos extremos, primeiro e quarto, respectivamente. 4.1.2 Componente de Conhecimento Específico Ao responderem à Questão 2 – “ Qual o grau de dificuldade desta prova na parte de Componente Específico? ” – 50,9 \% do grupo de estudantes classificaram -na como difícil ou muito difícil.
Além disso, o Componente de Conhecimento Específico da prova foi considerado com grau de dificuldade médio por 43,9 \% dos alunos ( Gráfico 4.3, Gráfico 4.4, e, no Anexo II, a Tabela II.2 ).
A análise das respostas dos estudantes quanto ao grau de dificuldade do Componente de Conhecimento Específico da prova, agregado por Grande Região, mostra que as diferenças entre a maior e as outras proporções de alunos que a avaliaram como difícil ou muito difícil são estatisticamente significativa.
A maior proporção foi a da região Sudeste ( 54,5 \% ) e as outras foram : 36,4 \% na região Norte, 37,8 \% na região Nordeste, 40,9 \% na Sul e 42,2 \% na Centro-Oeste.
O percentual de alunos que classificaram o grau de dificuldade como médio, no Componente de Conhecimento Específico, variou de 41,0 \% a 59,6 \%, para as regiões Sudeste e Norte, respectivamente.
Considerando -se a avaliação da dificuldade das questões do Componente de Conhecimento Específico da prova, de acordo com o desempenho dos estudantes, observa -se que há diferenças estatisticamente significativas do resultado do quarto de melhor desempenho em relação a todos os demais.
A proporção dos que classificaram a parte específica como difícil ou muito difícil variou de 32,6 \% ( 4º quarto ) a 60,3 \% ( 2º quarto ).
A alternativa modal para a Questão 2 foi o grau médio, com 38,6 \% do quarto inferior e 56,8 \% do superior optando por esta resposta. 4.2 EXTENSÃO DA PROVA EM RELAÇÃO AO TEMPO TOTAL Indagados quanto à extensão da prova, em relação ao tempo total oferecido para a sua resolução ( Questão 3 ), os estudantes apontaram, com maior incidência, a alternativa que considerava a extensão adequada, para todas as agregações consideradas ( Gráfico 4.5, Gráfico 4.6, e, no Anexo II, a Tabela II.3 ).
O percentual de alunos que responderam ser a extensão da prova adequada foi de 58,9 \%.
Já 32,5 \% dos inscritos presentes consideraram que o exame foi longo ou muito longo e menos do que 8,6 \% o avaliaram como curto ou muito curto.
Entre as Grandes Regiões a proporção daqueles que avaliaram a prova como longa ou muito longa em relação ao tempo total destinado à sua resolução variou pouco : de 30,3 \% na região Centro-Oeste até 37,9 \% na região Sul.
As diferenças entre as regiões não são estatisticamente significativas.
Considerando -se o desempenho dos alunos, nota -se ainda que 60,8 \% consideraram a extensão da prova adequada no quarto de desempenho inferior e 60,6 \% no de melhor desempenho ( quarto superior ).
Nos quartos intermediários, esta proporção foi 58,6 \% no segundo quarto e 55,5 \% no terceiro.
No Gráfico 4.6, pode -se constatar que há uma tendência crescente da proporção de estudantes que consideraram a prova longa ou muito longa em função dos quartos de desempenho nos três primeiros quartos.
Observa -se que não há diferença estatisticamente significativa entre as proporções de escolha desta resposta no primeiro e no último quarto. 4.3 COMPREENSÃO DOS ENUNCIADOS DAS QUESTÕES 4.3.1 Componente de Formação Geral Com relação aos enunciados das questões do Componente de Formação Geral ( Questão 4 ), as opiniões foram positivas, já que 78,3 \% dos alunos avaliados consideraram os enunciados de todas ou da maioria das questões claros e objetivos ( Gráfico 4.7, Gráfico 4.8, e, no Anexo II, a Tabela II.4 ).
Na análise regional, a percentagem de estudantes que avaliaram que todos ou a maioria dos enunciados das questões do Componente de Formação Geral estavam claros e objetivos variou de 75,8 \% na região Norte a 85,3 \% na região Centro-Oeste, sendo esta diferença não estatisticamente significativa.
A análise das percepções dos estudantes sobre a clareza e objetividade dos enunciados permite afirmar que todos, ou a maioria dos enunciados de questões relativas ao Componente de Formação Geral, foram considerados claros e objetivos para a maior parte dos respondentes ( maior do que 75 \% em todas as regiões e maior do que 70 \% para todos os quartos de desempenho ).
Segundo o desempenho, observa -se que a proporção dos que emitiram esta opinião cresce conforme o desempenho aumenta, até o 3º quarto, com diferença estatisticamente significativa do primeiro para o segundo quarto.
No quarto superior, a clareza e objetividade de todos ou da maioria dos enunciados das questões foi percebida por 80,1 \% dos alunos sem que seja estatisticamente significativa a diferença deste para o quarto imediatamente inferior, onde a proporção foi de 81,5 \%. 4.3.2 Componente de Conhecimento Específico Com relação aos enunciados das questões do Componente de Conhecimento Específico da prova, para 79,8 \% dos estudantes avaliados da Área de Engenharia de Controle e Automação, a clareza e a objetividade ( Questão 5 ) estavam presentes em todas ou na maioria das questões ( Gráfico 4.9, Gráfico 4.10, e no Anexo II, a Tabela II.5 ).
A maioria dos estudantes de todas as Grandes Regiões brasileiras considerou claros e objetivos todos ou a maioria dos enunciados das questões do Componente de Conhecimento Específico da prova, percentual sempre maior ou igual a 76,8 \%.
A diferença entre a maior e a menor proporção, 82,6 \% ( Sul ) e 76,8 \% ( Norte ), não é estatisticamente significativa.
A proporção de estudantes que consideraram os enunciados das questões claros e objetivos apresenta uma tendência crescente em relação ao aumento de desempenho : mais elevada no quarto superior ( 87,3 \% ) se comparada ao quarto inferior de desempenho ( 69,8 \% ).
As diferenças do primeiro para os demais quartos e do segundo para os dois seguintes são estatisticamente significativas. 4.4 SUFICIÊNCIA DAS INFORMAÇÕES/INSTRUÇÕES FORNECIDAS Ao avaliarem as informações/instruções fornecidas para a resolução das questões ( Questão 6 ), 86,8 \% dos respondentes da Área de Engenharia de Controle e Automação de todo o Brasil afirmaram que estas eram até excessivas ou suficientes em todas ou na maioria das questões ( Gráfico 4.11, Gráfico 4.12, e, no Anexo II, a Tabela II.6 ).
Quanto à distribuição de respondentes pelas Grandes Regiões observa -se que a proporção de estudantes que consideraram as informações/instruções fornecidas até excessivas ou suficientes em todas ou na maioria das questões foi sempre maior ou igual a 85,7 \%, chegando a 92,4 \% na região Sul.
A diferença entre o maior percentual, obtido na região Sul ( 92,4 \% ), e o menor, da região Sudeste ( 85,7 \% ) é estatisticamente significativa.
Levando -se em conta o desempenho dos participantes, nota -se diferença estatisticamente significativa entre as opiniões de estudantes do quarto inferior e superior de desempenho, como mostra o Gráfico 4.12.
O percentual de participantes que avaliaram as informações/instruções como até excessivas ou suficientes em todas ou na maioria das questões foi mais elevado no quarto superior ( 94,4 \% ), percentual superior à média nacional ( 86,8 \% ).
Já no quarto inferior, a suficiência das informações/instruções declarada como até excessiva em todas ou na maioria das questões foi percebida por 79,5 \% dos respondentes.
Estas proporções são crescentes com os quartos de desempenho. 4.5 DIFICULDADE ENCONTRADA AO RESPONDER À PROVA Perguntados sobre as dificuldades com as quais se depararam ao responder à prova ( Questão 7 ), 16,0 \% dos estudantes apontaram o desconhecimento do conteúdo.
Para 47,6 \%, a forma diferente de abordagem do conteúdo foi indicada como dificuldade.
Já a falta de motivação para fazer a prova foi a dificuldade apontada por 15,5 \% dos respondentes.
Considerando -se todo o Brasil, 15,9 \% dos respondentes afirmaram que não tiveram qualquer tipo de dificuldade para responder à prova ( Tabela II.7 no Anexo II ).
Os Gráficos 4.13 e 4.14 apresentam os percentuais de estudantes que apontaram o desconhecimento do conteúdo como dificuldade percebida ao responder à prova.
Na análise por Grandes Regiões, o percentual de inscritos e presentes que apontaram o desconhecimento do conteúdo como dificuldade ao responder à prova não superou 20,2 \%.
Os percentuais variaram de 14,9 \% na região Sudeste a 20,2 \% na Norte, sendo as diferenças entre as regiões não significativas estatisticamente.
A forma diferente de abordagem do conteúdo foi a escolha modal dos estudantes, com percentuais que variaram de 37,3 \% ( região Sul ) a 50,1 \% ( Sudeste ).
O percentual de alunos que citaram a falta de motivação como dificuldade variou de 13,6 \% ( região Sudeste ) a 23,9 \% ( região Centro-Oeste ).
Os que declararam não ter qualquer dificuldade para responder à prova variaram de 9,2 \% na região Centro-Oeste a 18,0 \% na Sul.
Com relação aos quartos de desempenho, o desconhecimento do conteúdo foi a opção escolhida por 17,1 \% dos estudantes do quarto superior e 15,0 \% do quarto inferior.
A alternativa modal para os alunos, quando agregados pelos quartos de desempenho, foi que a dificuldade encontrada foi causada pela forma diferente de 98 abordagem do conteúdo : 51,6 \% no quarto inferior e 36,5 \% do quarto superior assim o responderam. 4.6 CONTEÚDOS DAS QUESTÕES OBJETIVAS DA PROVA Ao analisarem os conteúdos das questões objetivas da prova ( Questão 8 ), um percentual muito pequeno dos estudantes avaliados, apenas 3,1 \%, afirmou que não estudou ainda a maioria desses conteúdos ( Gráficos 4.15, Gráfico 4.16, e a Tabelas II.8 no Anexo II ).
A maioria ( 68,3 \% ) afirmou ter estudado e aprendido muitos ou todos os conteúdos.
Na análise por Grande Região, a proporção de respondentes que escolheram a opção não estudou ainda a maioria desses conteúdos foi pequena.
Observa -se que nas regiões Norte ( 4,0 \% ) e Nordeste ( 5,1 \% ), apesar de pequenas, as proporções foram maiores do que a média nacional ( 3,1 \% ).
Observa -se que entre as regiões Norte e Nordeste não há diferença estatisticamente significativa.
Em todas as regiões, a maioria dos presentes afirmou ter estudado e aprendido muitos ou todos os conteúdos, com proporções variando entre 64,5 \% na região Nordeste e 79,2 \% na Sul.
Considerando -se separadamente as opiniões de estudantes dos quatro quartos de desempenho, observa -se que, no quarto inferior, 5,6 \% ofereceram como resposta que não estudou ainda a maioria desses conteúdos, sendo 1,3 \% os do quarto superior com a mesma resposta.
A diferença entre os alunos que optaram por este motivo de dificuldade nos quartos extremos é estatisticamente significativa.
Tendo em conta o quarto superior, 85,8 \% dos alunos afirmaram ter estudado e aprendido muitos ou todos os conteúdos.
No outro extremo, no primeiro quarto, 49,7 \% optaram pelas mesmas categorias. 4.7 TEMPO GASTO PARA CONCLUIR A PROVA Ao responderem sobre o tempo de conclusão da prova ( Questão 9 ), mais da metade dos estudantes ( 74,2 \% ) afirmou ter gasto entre duas e quatro horas ( Gráfico 4.17, Gráfico 4.18 e, no Anexo II, a Tabela II.9 ).
Considerando -se as cinco Grandes Regiões brasileiras, o percentual dos que utilizaram entre duas e quatro horas para finalizar a prova apenas nas regiões Nordeste ( 67,6 \% ) e Centro-Oeste ( 70,1 \% ) foi inferior ao percentual nacional.
Nas demais Grandes Regiões o percentual de alunos que dispensaram entre duas e quatro horas para concluir a prova, foi superior ao percentual nacional, como mostra o Gráfico 4.17.
Uma vez sendo analisadas as alternativas escolhidas pelos estudantes que se situam nos diferentes quartos de desempenho, observa -se uma tendência crescente : uma maior proporção de participantes no quarto superior declarou ter gasto entre duas e quatro horas para concluir a prova quando comparadas com os dos quartos inferiores.
As diferenças entre o primeiro ( 66,2 \% ) e os demais quartos são estatisticamente significativas, conforme mostra o Gráfico 4.18.

% matched lemmas: 
\end{textsample}
